% ============================================================
%  ClearDrive — Hybrid AI and Image Processing Dashboard
%  for Low-Visibility Driving
%
%  DIP Course Project Report
%  IIITDM Kancheepuram — CSE Department
%
%  Team: Gyan Chandra (CS23I1053)
%        Vishal Singh  (CS23I1039)
%        Ranveer Gautam (ME23B2031)
%
%  Compile on Overleaf: single pdflatex run is sufficient.
% ============================================================

\documentclass[12pt, a4paper]{article}

% ── Page Layout ────────────────────────────────────────────
\usepackage[
  top=1in, bottom=1in,
  left=1.1in, right=1.1in
]{geometry}

% ── Fonts & Encoding ───────────────────────────────────────
\usepackage[T1]{fontenc}
\usepackage[utf8]{inputenc}
\usepackage{lmodern}
\usepackage{microtype}

% ── Mathematics ────────────────────────────────────────────
\usepackage{amsmath}
\usepackage{amssymb}

% ── Graphics ───────────────────────────────────────────────
\usepackage{graphicx}
\usepackage{float}

% ── Tables ─────────────────────────────────────────────────
\usepackage{booktabs}
\usepackage{array}

% ── Colors & Boxes ─────────────────────────────────────────
\usepackage{xcolor}
\usepackage{tcolorbox}
\tcbuselibrary{skins, breakable}

% ── Section Styling ────────────────────────────────────────
\usepackage{titlesec}
\titleformat{\section}
  {\large\bfseries\color{blue!60!black}}
  {\thesection.}{0.6em}{}[\titlerule]
\titleformat{\subsection}
  {\normalsize\bfseries\color{blue!40!black}}
  {\thesubsection}{0.5em}{}

% ── Code Listings ──────────────────────────────────────────
\usepackage{listings}
\lstset{
  basicstyle=\ttfamily\small,
  backgroundcolor=\color{gray!10},
  frame=single,
  rulecolor=\color{gray!60},
  breaklines=true,
  keywordstyle=\color{blue!70!black}\bfseries,
  commentstyle=\color{green!50!black}\itshape,
  stringstyle=\color{red!60!black},
  showstringspaces=false,
  numbers=left,
  numberstyle=\tiny\color{gray},
  numbersep=6pt,
  language=Python
}

% ── Header / Footer ────────────────────────────────────────
\usepackage{fancyhdr}
\pagestyle{fancy}
\fancyhf{}
\fancyhead[L]{\small\textit{ClearDrive — DIP Course Project}}
\fancyhead[R]{\small\textit{IIITDM Kancheepuram}}
\fancyfoot[C]{\thepage}
\renewcommand{\headrulewidth}{0.4pt}

% ── Hyperlinks ─────────────────────────────────────────────
\usepackage{hyperref}
\hypersetup{
  colorlinks=true,
  linkcolor=blue!70!black,
  citecolor=blue!70!black,
  urlcolor=teal,
  pdftitle={ClearDrive Project Report},
  pdfauthor={Gyan Chandra, Vishal Singh, Ranveer Gautam}
}

% ── Line Spacing ───────────────────────────────────────────
\usepackage{setspace}
\setstretch{1.35}

% ── Caption Styling ────────────────────────────────────────
\usepackage{caption}
\captionsetup{
  font=small,
  labelfont=bf,
  skip=6pt
}

% ── Paragraph Spacing ──────────────────────────────────────
\setlength{\parskip}{6pt}
\setlength{\parindent}{0pt}

% ============================================================
%   TITLE PAGE (FIXED FOR EXACT 1 PAGE FIT)
% ============================================================
\begin{document}

\begin{titlepage}
  % Temporarily reduce line spacing just for the title page to prevent overflow
  \begin{spacing}{1.1}
  \centering
  
  % Negative vspace pulls everything slightly up
  \vspace*{-1.0cm}

  %% ── Institute Logo ──
  \includegraphics[width=4cm, height=4cm, keepaspectratio]{iiit_logo.png}

  \vspace{0.4cm}
  {\LARGE \textbf{Indian Institute of Information Technology,}\\[5pt]
           \textbf{Design and Manufacturing (IIITDM), Kancheepuram}}

  \vspace{0.2cm}
  {\large Department of Computer Science and Engineering}

  \vspace{0.6cm}
  \textcolor{blue!60!black}{\rule{\linewidth}{2pt}}
  \vspace{0.3cm}

  {\Huge \textbf{ClearDrive}}\\[8pt]
  {\Large Hybrid AI and Image Processing Dashboard\\
          for Low-Visibility Driving}

  \vspace{0.3cm}
  \textcolor{blue!60!black}{\rule{\linewidth}{2pt}}

  \vspace{0.6cm}
  {\large \textbf{Digital Image Processing (DIP)\\Course Project Report}}

  \vspace{0.6cm}
  \begin{tcolorbox}[
    colback=blue!5!white,
    colframe=blue!50!black,
    width=0.85\linewidth,
    arc=4pt
  ]
    \centering
    \begin{tabular}{lll}
      \toprule
      \textbf{Name} & \textbf{Roll No.} & \textbf{Programme}\\
      \midrule
      Gyan Chandra   & CS23I1053 & Dual Degree CSE\\
      Vishal Singh   & CS23I1039 & Dual Degree CSE\\
      Ranveer Gautam & ME23B2031 & B.Tech Smart Manufacturing\\
      \bottomrule
    \end{tabular}
  \end{tcolorbox}

  \vfill
  {\large \textbf{April 2026}}
  
  \end{spacing}
\end{titlepage}

% ── Table of Contents ──────────────────────────────────────
\tableofcontents
\newpage

% ============================================================
\section{Introduction}
% ============================================================

\subsection{The Problem: Driving in Bad Weather in India}

Anyone who has driven on Indian highways during early-morning winter fog,
or through monsoon mist on a hill-road, knows how dangerous it can get.
When fog is heavy, you cannot see lane markings. You cannot tell if there is
a pedestrian or a stalled truck twenty meters ahead. Night driving on
unlit state highways is equally risky — headlights barely penetrate the
darkness, and the contrast between the road and its surroundings becomes
almost zero.

This is not just a comfort issue — it is a proven safety crisis. A significant
fraction of road accidents in India occur in low-visibility conditions, and a
large number of those happen on national highways. The root cause is simple:
the human eye, and by extension, the camera-based systems inside modern cars,
cannot \textit{see} well when fog scatters light or when there is very little
ambient light at night.

\subsection{Why We Built a Camera-Based System (Not LiDAR)}

Modern high-end cars use LiDAR (Light Detection and Ranging) sensors to build
a 3D map of the surroundings. LiDAR is accurate and works in the dark, but it
has two major problems for real-world Indian deployment:

\begin{itemize}
  \item \textbf{Cost:} A single automotive-grade LiDAR unit costs anywhere from
  ₹3~lakhs to ₹30~lakhs. This is simply not practical for the average Indian car buyer.

  \item \textbf{No semantic understanding:} LiDAR gives you a 3D point cloud.
  It tells you \textit{where} objects are, but not \textit{what} they are. It
  cannot read lane markings painted on the road because those markings are flat
  and have the same reflectivity as the surrounding tarmac. A camera, on the
  other hand, sees \textit{colour and texture} — exactly the information that
  lane markings, traffic signs, and road surfaces carry.
\end{itemize}

Our solution, \textbf{ClearDrive}, uses only a standard dashcam (a plain
RGB camera) combined with image processing algorithms and pre-trained AI
models running on an ordinary laptop GPU (NVIDIA RTX 3050). The total hardware
cost is effectively zero if a car already has a dashcam. All the intelligence
is in the software.

\subsection{What ClearDrive Does}

In simple terms, ClearDrive watches the video coming from a dashcam in real
time, frame by frame, and does the following:

\begin{enumerate}
  \item It first \textbf{identifies what kind of bad weather} is present
  (foggy, night, or clear) using an AI model.
  \item Based on that, it \textbf{cleans up the image} — either by removing
  the haze mathematically (fog mode) or by boosting the brightness intelligently
  (night mode).
  \item It then \textbf{finds the lane markings} using classical edge detection.
  \item As a backup when lines are invisible, it uses a \textbf{deep learning
  model} to highlight the drivable road area in green.
  \item It shows all of this on a \textbf{4-panel live dashboard} and logs
  performance data to a file.
\end{enumerate}

% ============================================================
\section{System Flow and Architecture}
% ============================================================

\subsection{The Big Picture: From Standalone Pipeline to Client-Server Architecture}

The original ClearDrive pipeline has been extended into a full-stack,
client-server system. The DIP processing core now runs as a web service inside
a \textbf{FastAPI} backend, while a browser-based \textbf{React} dashboard
acts as the client. Figure~\ref{fig:flowchart} illustrates the complete
end-to-end data flow.

\begin{figure}[H]
  \centering
  \includegraphics[width=0.85\textwidth]{flowchart.png}
  \caption{ClearDrive end-to-end data flow: raw video is ingested by the
           FastAPI backend, processed through the four-stage AI-DIP pipeline,
           and the resulting MJPEG stream and JSON metrics are consumed by
           the React frontend to render the live driver HUD.}
  \label{fig:flowchart}
\end{figure}

At a high level, the data flow is as follows:

\begin{enumerate}
  \item \textbf{Raw Video Input} — A dashcam video file, a live webcam feed,
  or a user-uploaded clip is passed to the backend.
  \item \textbf{Backend Processing} — The FastAPI server runs the full
  four-stage DIP pipeline (weather classification $\rightarrow$ dehazing/
  CLAHE $\rightarrow$ lane detection $\rightarrow$ HUD assembly) on each
  frame and encodes the composite dashboard as a JPEG.
  \item \textbf{Network Transport} — Processed frames are streamed to the
  browser over HTTP using the \textbf{MJPEG} (multipart/x-mixed-replace)
  protocol. Real-time metrics (FPS, Visibility Score, Contrast Gain, and
  weather mode) are served separately as a lightweight JSON payload via the
  \texttt{/metrics} REST endpoint, polled by the frontend every second.
  \item \textbf{Frontend Display} — The React dashboard renders the live
  MJPEG stream and overlays the numeric metrics in the HUD panel, giving the
  operator a real-time view of system state.
\end{enumerate}

\subsection{Backend API Endpoints}

The backend exposes four HTTP endpoints, each of which triggers the same
underlying \texttt{process\_and\_stream()} generator function:

\begin{center}
\begin{tabular}{lll}
  \toprule
  \textbf{Endpoint} & \textbf{Method} & \textbf{Purpose}\\
  \midrule
  \texttt{/health}             & GET  & Heartbeat — confirms backend is alive\\
  \texttt{/video\_feed}        & GET  & Streams the pre-recorded dashcam file\\
  \texttt{/webcam\_feed}       & GET  & Streams from connected camera (index~0)\\
  \texttt{/upload}            & POST & Accepts an uploaded video; returns a \texttt{stream\_id}\\
  \texttt{/stream/\{id\}}     & GET  & Streams previously uploaded file by ID\\
  \texttt{/metrics}           & GET  & Returns live JSON metrics snapshot\\
  \bottomrule
\end{tabular}
\end{center}

\subsection{Stage 1: AI Weather Classification (The Gatekeeper)}

The first thing the system does with every frame is ask: \textit{``What kind
of weather am I looking at?''} This is handled by the \texttt{WeatherClassifier}
class in \texttt{classifier.py}.

Internally, this uses \textbf{MobileNetV2}, a lightweight deep neural network
pre-trained on ImageNet. We replaced its final output layer to produce three
outputs instead of 1000: \textbf{CLEAR}, \textbf{FOGGY}, or \textbf{NIGHT}.

The input frame is resized to $224 \times 224$ pixels, normalised using standard
ImageNet mean and standard deviation values, and then fed into the model running
on the GPU. The class with the highest score is the predicted weather. To reduce
redundant GPU inference calls over the network, classification is run only once
every 10 frames; the cached result is reused for intermediate frames.

\subsection{Stage 2: Frame Processing Based on Weather}

Once we know the weather, we route the frame to the right processing arm:

\begin{tcolorbox}[
  colback=green!5!white,
  colframe=green!50!black,
  title={\textbf{Routing Logic (from backend/api.py)}},
  arc=4pt
]
\begin{lstlisting}
if weather_mode == "NIGHT":
    processed_img = enhance_low_light(frame)   # CLAHE

elif weather_mode == "FOGGY":
    # Dark Channel Prior + CLAHE reinforcement
    processed_img = recover_image(img_float, t_smoothed, A_smoothed, t0=0.05)

else:  # CLEAR
    processed_img = frame.copy()   # No change needed
\end{lstlisting}
\end{tcolorbox}

\subsection{Stage 3: Lane Detection (Edges + AI Backup)}

After the image is cleaned up, the system runs lane detection on the
processed frame:

\begin{itemize}
  \item \textbf{Canny Edge Detection} finds the sharp boundaries (like lane
  markings) in the image.
  \item \textbf{Hough Transform} connects those edge points into actual lines.
  \item If these lines are not found (because the fog was too thick even after
  dehazing), the \textbf{LR-ASPP MobileNetV3} segmentation model kicks in
  and highlights the drivable road area in green.
\end{itemize}

\subsection{Stage 4: HUD Assembly and Network Streaming}

After all processing stages complete for a given frame, the backend:
\begin{itemize}
  \item Generates a coloured \textbf{Visibility Heat Map} from the transmission
  data.
  \item Calculates \textbf{Contrast Gain} and \textbf{Visibility Score} metrics
  and writes them into the shared \texttt{system\_status} dictionary.
  \item Assembles the \textbf{4-panel dashboard} (RAW INPUT, AI-ENHANCED,
  HUD, AI SEGMENT) into a single $1280 \times 720$ frame.
  \item JPEG-encodes the frame and yields it as an MJPEG chunk over the open
  HTTP connection to the browser client.
\end{itemize}

% ============================================================
\section{Practical Implementation Details}
% ============================================================

\subsection{Software Stack Overview}

The ClearDrive system is built on a two-tier software stack that keeps the DIP
processing logic strictly on the server, while the browser client handles only
presentation.

\textbf{Backend} (\texttt{backend/} directory): Written in \textbf{Python 3.10}
and served using \textbf{FastAPI} with \textbf{Uvicorn} as the ASGI server.
All image processing runs through OpenCV (v4.x) and PyTorch (v2.x). The backend
exposes MJPEG streams and a JSON metrics endpoint; it handles CORS so the React
client can connect from any origin during development.

\textbf{Frontend} (\texttt{frontend/} directory): A single-page application
built with \textbf{React 19} and bundled with \textbf{Vite}. The dashboard
renders the live MJPEG stream inside an \texttt{<img>} element (the browser's
native MJPEG decoding is used, requiring no additional video codec library) and
polls the \texttt{/metrics} endpoint on a one-second interval to update the
HUD overlays — FPS counter, Visibility Score, Contrast Gain, and the active
weather mode badge.

This separation ensures that the computationally expensive DIP pipeline remains
decoupled from the display layer. Swapping the frontend framework or deploying
the backend on a more powerful edge server would require no changes to the core
image processing code.

This section explains the most interesting parts of our implementation — the
problems we faced and how we solved them.

\subsection{Problem 1: The AI Was Flickering Between Modes}

\textbf{The Issue:} When we first built the system, we called the weather
classifier on every single frame independently. The result was a mess. The AI
would say ``FOGGY'' for 8 frames, then ``CLEAR'' for 1 frame (because a truck
passed and changed the brightness), then ``FOGGY'' again. This made the
processed image flash between the dehazed version and the raw version, which
was visually irritating and practically useless.

\textbf{Our Solution — Temporal Smoothing with a 10-Frame Buffer:}

We created a simple Python list called \texttt{mode\_history} that stores the
weather predictions from the last 10 frames. Instead of using the raw prediction,
we take a \textit{majority vote} over this list — whichever class appears most
often in the last 10 frames is the one we actually use.

\begin{tcolorbox}[
  colback=blue!5!white,
  colframe=blue!50!black,
  title={\textbf{Temporal Smoothing Logic (main.py, lines 56--61)}},
  arc=4pt
]
\begin{lstlisting}
mode_history.append(raw_mode)

# Keep only the last 10 frames in memory
if len(mode_history) > 10:
    mode_history.pop(0)

# The final decision: pick the most frequent class
weather_mode = max(set(mode_history), key=mode_history.count)
\end{lstlisting}
\end{tcolorbox}

Think of it like a polling system: 10 voters (frames) each cast their vote for
a weather mode. The mode that gets the most votes wins. A single rogue vote
(one bad frame) cannot change the outcome. This simple trick completely
eliminated the flickering problem.

We applied the same idea to the \textbf{Atmospheric Light} ($A$) estimated in
the DCP algorithm, using a 5-frame average. This stopped sudden brightness jumps
between consecutive foggy frames.

\subsection{The Main DIP Algorithm: Dark Channel Prior (DCP) for Fog Removal}
\label{sec:dcp}

This is the heart of the fog removal system, implemented in \texttt{dehaze.py}.
The algorithm is based on a clever observation about natural images: in almost
every small patch of an outdoor image, at least one colour channel in at least
one pixel will be very dark (close to zero intensity). This is called the
\textbf{Dark Channel}. When fog is present, this dark channel is no longer dark
— the haze veil artificially brightens it.

The entire math of fog can be written as one equation (the \textbf{Koschmieder
Haze Model}):

\begin{equation}
  I(x) = J(x) \cdot t(x) + A \cdot (1 - t(x))
  \label{eq:haze}
\end{equation}

Where:
\begin{itemize}
  \item $I(x)$ is what the camera actually captures (the hazy image).
  \item $J(x)$ is the original fog-free scene — \textbf{this is what we want to recover}.
  \item $t(x)$ is the \textbf{transmission map}: a value between 0 and 1 that
  tells us what fraction of the original scene light actually reached the camera.
  A value close to 1 means clear; close to 0 means heavy fog.
  \item $A$ is the \textbf{atmospheric light}: the overall brightness/colour of
  the fog itself (the grey-white veil).
\end{itemize}

If we can estimate $t(x)$ and $A$, we can rearrange Equation~\ref{eq:haze} to
get $J(x)$, the dehazed image.

\subsubsection{Step 1: Calculate the Dark Channel}

For each pixel, we take the \textit{minimum} across all three colour channels
(R, G, B). Then we take the \textit{minimum} of that within a $15 \times 15$
pixel neighbourhood around it. This gives us a single 2D map — the dark channel:

\begin{equation}
  J^{\text{dark}}(x) =
  \min_{c \in \{R,G,B\}}
  \Bigl(
    \min_{y \in \text{patch}(x)} I^c(y)
  \Bigr)
  \label{eq:dark_channel}
\end{equation}

In code, the neighbourhood minimum is done efficiently using
\texttt{cv2.erode()} with a $15 \times 15$ rectangular kernel.

\subsubsection{Step 2: Estimate Atmospheric Light ($A$)}

Fog is bright. So the pixels in the dark channel that are still relatively bright
are the ones most dominated by fog. We pick the \textbf{top 0.1\% brightest}
pixels of the dark channel, look up their actual RGB colours in the original
image, and take the average. That average colour is our estimate of $A$.

\subsubsection{Step 3: Estimate the Transmission Map ($t$)}

The transmission at each pixel tells us how much scene radiance got through the
fog. We calculate it as:

\begin{equation}
  t(x) = 1 - \omega \cdot \min_{c}
  \Bigl(
    \min_{y \in \text{patch}(x)} \frac{I^c(y)}{A^c}
  \Bigr)
  \label{eq:transmission}
\end{equation}

The key parameter here is $\omega$ (omega). A value of $\omega = 1.0$ would
remove 100\% of the haze. We used $\omega = \mathbf{0.95}$ — very aggressive.
We chose this higher value deliberately because for a driving safety system,
we would rather over-correct the fog and have a slightly over-saturated image
than under-correct it and leave the driver with a still-foggy view.

We used a smaller $7 \times 7$ patch size here (compared to $15 \times 15$ for
the dark channel) to capture finer spatial structure in the transmission map.
The raw transmission map can be blocky, so we smooth it with a $5 \times 5$
Gaussian blur.

\subsubsection{Step 4: Recover the Clean Image}

Now we just rearrange Equation~\ref{eq:haze}:

\begin{equation}
  J(x) = \frac{I(x) - A}{\max(t(x),\; t_0)} + A
  \label{eq:recover}
\end{equation}

We clip $t(x)$ from below at $t_0 = 0.10$ to prevent division-by-very-small-number
in extremely foggy pixels. The result $J(x)$ is our dehazed frame.

\subsection{Night Mode: CLAHE Enhancement}

For night frames, fog removal is the wrong tool — the problem is not scattering
but low light. We use \textbf{CLAHE (Contrast-Limited Adaptive Histogram
Equalization)}, implemented in \texttt{enhancement.py}.

The trick is to work in the \textbf{CIE-LAB colour space} instead of RGB.
LAB separates the image into:
\begin{itemize}
  \item \textbf{L} — Lightness (how bright a pixel is)
  \item \textbf{a, b} — Colour information
\end{itemize}

We apply CLAHE only to the \textbf{L channel}. This boosts the contrast and
brightness without changing the colours. If we applied it to the RGB channels
directly, the image would become unnaturally saturated (too colourful).

Our CLAHE settings: \texttt{clipLimit = 3.0} (limits noise amplification) and
\texttt{tileGridSize = (8, 8)} (the image is divided into $8 \times 8$ local
tiles, each of which gets its own histogram equalization). After processing, we
merge the enhanced L channel back with the original a and b channels and convert
back to BGR.

\subsection{Lane Detection: Canny + Hough Transform}

Lane detection runs in \texttt{lanes.py} and follows a classic 4-step pipeline:

\textbf{Step 1 — Greyscale and Blur:} Convert the (now cleaned-up) frame to greyscale,
then apply a $5 \times 5$ Gaussian blur to reduce noise. Noise in the image creates
many false edges, so blurring first makes the Canny algorithm much more reliable.

\textbf{Step 2 — Canny Edge Detection:} The Canny algorithm finds pixels where
the image brightness changes sharply — i.e., edges. We use two thresholds:
\begin{itemize}
  \item \textbf{Low threshold = 50:} Pixel is a potential edge.
  \item \textbf{High threshold = 150:} Pixel is definitely an edge. Any pixel
  connected to a strong edge that also crosses 50 is kept; everything else
  is discarded.
\end{itemize}

\textbf{Step 3 — Region of Interest Masking:} We don't care about edges everywhere
in the image. We only care about the road ahead. We draw a triangle/trapezoid in
the lower half of the frame and discard all edges outside this region using a
binary mask and a bitwise AND operation.

\textbf{Step 4 — Probabilistic Hough Transform:} The Hough Transform takes the
scattered edge pixels and assembles them into line segments. Each point in image
space maps to a sinusoidal curve in a parameter space (called Hough Space). Where
many curves intersect, a strong line exists in the image. Our parameters:
minimum vote count = 100, minimum line length = 40 pixels, maximum gap = 5 pixels.

\subsection{The Deep Learning Backup: Semantic Road Segmentation}

\textbf{The Problem it Solves:} Even after DCP dehazing, if the fog was
extremely dense, the lane markings might still be invisible. In this case,
the Hough Transform returns zero lines, and we trigger a ``LANE LOST'' warning.
But the system still needs to show the driver \textit{where the road is}.

\textbf{Our Solution} (in \texttt{segmentation.py}): We use a pre-trained deep
learning model called \textbf{LR-ASPP MobileNetV3-Large}. This model was trained
on the COCO dataset and can classify every pixel in an image into 21 categories
(road, car, person, sky, etc.). We use it to find all pixels belonging to the
``road'' or ``background flat surface'' classes (class IDs 0 and 7 in the COCO
dataset) and colour them green, creating a road-surface highlight.

\begin{figure}[H]
  \centering
  \includegraphics[width=0.85\textwidth]{segmentation_focus.png}
  \caption{The LR-ASPP MobileNetV3 model highlighting the drivable road area
           in green, even when lane lines are not visible.}
  \label{fig:seg}
\end{figure}

\textbf{The Ego-Vehicle Masking Trick:} When we first ran the segmentation,
we noticed two problems:

\begin{enumerate}
  \item The sky was sometimes classified as ``background'' and getting coloured green.
  \item The bottom 15\% of the frame shows the car's bonnet or wiper blades,
  which confused the model.
\end{enumerate}

Our fix was beautifully simple: we \textbf{zero out the top 60\% of the
segmentation mask}. This means we only display the green road highlight in the
bottom 40\% of the frame — exactly where the road is in a typical dashcam
view. The sky false-positives disappear, and the bonnet area is naturally
excluded from the useful region. This was done with one line of code:

\begin{lstlisting}
road_mask[0 : int(height * 0.6), :] = 0  # Blank out top 60%
\end{lstlisting}

The final road mask is then blended over the processed frame at 30\% transparency
using \texttt{cv2.addWeighted()}, giving a glass-like green overlay on the road.

% ============================================================
\section{Dashboard and Results}
% ============================================================

\subsection{The 4-Panel Dashboard}

All the processed outputs are combined into a single $1280 \times 720$ HD video
frame for display and recording. The dashboard has four quadrants (see
Figure~\ref{fig:dashboard}):

\begin{figure}[H]
  \centering
  \includegraphics[width=\textwidth]{dashboard_full.png}
  \caption{The ClearDrive 4-panel real-time dashboard as seen in the final output video.}
  \label{fig:dashboard}
\end{figure}

\begin{center}
\begin{tabular}{cll}
  \toprule
  \textbf{Panel} & \textbf{Label in Dashboard} & \textbf{What It Shows}\\
  \midrule
  Top-Left     & RAW INPUT FEED          & Original camera footage + alert text\\
  Top-Right    & AI-ENHANCED STREAM      & The fog-removed or night-boosted image\\
  Bottom-Left  & RELIABILITY HUD         & Colour-coded visibility heat map\\
  Bottom-Right & AI ROAD SEGMENTATION    & Green road highlight from MobileNetV3\\
  \bottomrule
\end{tabular}
\end{center}

A \textbf{header bar} runs across the top of the entire dashboard. It shows the
current mode (FOGGY, NIGHT, or CLEAR), live FPS, Visibility Score, and Contrast
Gain. The header turns \textit{red} when an alert is active and \textit{green}
when everything is safe.

\subsection{The Visibility Heat Map (HUD)}

The Reliability HUD (bottom-left panel) is generated from the transmission map
$t(x)$ that the DCP algorithm computes. We colour-code every pixel based on
how much scene radiance passed through the fog:

\begin{equation}
  \text{Colour}(x) = \begin{cases}
    \text{GREEN}  & \text{if } t(x) \geq 0.65 \quad (\text{safe zone})\\
    \text{YELLOW} & \text{if } 0.35 \leq t(x) < 0.65 \quad (\text{caution zone})\\
    \text{RED}    & \text{if } t(x) < 0.35 \quad (\text{danger zone})
  \end{cases}
\end{equation}

This coloured map is then heavily blurred with a $51 \times 51$ Gaussian kernel
so it looks like a smooth heatmap rather than a pixelated grid. It is blended
onto the original frame at 35\% transparency using \texttt{apply\_hud()}.

\subsection{Alert System}

The \texttt{alerts.py} module generates on-screen text warnings based on two
conditions:

\begin{itemize}
  \item If the \textbf{Visibility Score drops below 60\%}: show a
  \textcolor{orange}{\textbf{CAUTION}} warning. If it drops below 40\%:
  show a \textcolor{red}{\textbf{CRITICAL}} alert.
  \item If \textbf{zero lane lines are detected} by the Hough Transform: show
  a \textcolor{orange}{\textbf{WARNING: LANE LOST}} message.
\end{itemize}

These text messages are printed directly on the RAW INPUT quadrant in
colour-coded text so the driver sees an alert even while looking at the
original unprocessed view.

\subsection{Performance Metrics and Data Logging}

We measure two key numbers for every processed frame:

\subsubsection{Contrast Gain}

Contrast is measured as the \textbf{standard deviation of pixel intensities}
in the greyscale image. A foggy image has low standard deviation (everything
is the same grey-white colour). A clear image has high standard deviation
(bright road markings, dark shadows, vivid colours).

\begin{equation}
  \text{Contrast Gain} (\%) =
  \frac{\sigma(\text{dehazed}) - \sigma(\text{original})}{\sigma(\text{original}) + \epsilon}
  \times 100
\end{equation}

A positive Contrast Gain means our system successfully improved the visual
quality of the image. In our experiments on the foggy dashcam footage, we
consistently measured gains well above 40\%.

\subsubsection{Visibility Score}

This answers the question: \textit{``What percentage of the road ahead is
clearly visible?''} We count the fraction of pixels where the transmission
value $t(x) \geq 0.65$ (i.e., in the safe zone):

\begin{equation}
  V_s (\%) = \frac{\text{Number of pixels where } t(x) \geq 0.65}{\text{Total pixels}} \times 100
\end{equation}

\subsubsection{CSV Logging}

The \texttt{PerformanceLogger} class in \texttt{logger.py} opens a CSV file
called \texttt{performance\_results.csv} at startup and writes one row for every
processed frame. Each row has: \texttt{Timestamp, Mode, FPS, Contrast\_Gain,
Visibility\_Score}.

After the video is fully processed, we run \texttt{plot\_metrics.py} to generate
the analytics report shown in Figure~\ref{fig:metrics}.

\begin{figure}[H]
  \centering
  \includegraphics[width=\textwidth]{metrics_report.png}
  \caption{Post-processing analytics report generated from the logged CSV data,
           showing FPS stability, contrast gain over time, and visibility score.}
  \label{fig:metrics}
\end{figure}

The analytics report has three graphs:
\begin{enumerate}
  \item \textbf{FPS over time} — confirms the system ran at or above 20 FPS
  (real-time threshold) on the RTX 3050 GPU throughout the entire video.
  \item \textbf{Contrast Gain (\%) over frames} — shows how much the DCP
  algorithm improved each frame. Higher spikes correspond to the frames with
  the densest fog, where DCP had the most work to do.
  \item \textbf{Visibility Score (\%) over frames} — shows the evolution of
  safe-zone coverage throughout the video.
\end{enumerate}

% ============================================================
\section{Hardware and Software Setup}
% ============================================================

\begin{center}
\begin{tabular}{ll}
  \toprule
  \textbf{Component} & \textbf{Details}\\
  \midrule
  GPU          & NVIDIA GeForce RTX 3050 (4 GB VRAM)\\
  Framework    & Python 3.10, PyTorch 2.x, OpenCV 4.x\\
  AI Models    & MobileNetV2 (classification), LR-ASPP MobileNetV3-Large (segmentation)\\
  Weights      & Pre-trained on ImageNet and COCO 2017 (via \texttt{torchvision})\\
  Input        & \texttt{foggy\_dashcam.mp4} — real-world dashcam footage\\
  Processing   & Internal resolution: $640 \times 480$ px\\
  Output       & XVID AVI at $1280 \times 720$ px, 20 FPS\\
  Data Logging & CSV (pandas + matplotlib for offline analysis)\\
  \bottomrule
\end{tabular}
\end{center}

The processing resolution was intentionally kept at $640 \times 480$ to reduce
GPU load. All four modules (classifier, DCP, segmentor, Canny-Hough) run in
parallel on a single frame before the next frame is fetched, allowing the system
to maintain real-time throughput.

% ============================================================
\section{Conclusion}
% ============================================================

We built \textbf{ClearDrive}, a working real-time Driver Assistance system that
addresses one of the most common causes of road accidents in India: low-visibility
driving. The system runs on a standard camera and a mid-range laptop GPU, making
it practically deployable without exotic hardware.

The key achievement of this project is the seamless integration of two very
different technologies:

\begin{itemize}
  \item \textbf{Classical DIP algorithms} (DCP dehazing, CLAHE, Canny edges,
  Hough Transform) that are mathematically precise, interpretable, and fast.
  \item \textbf{Pre-trained Deep Learning models} (MobileNetV2 classifier,
  LR-ASPP MobileNetV3 segmentor) that bring semantic understanding — knowing
  \textit{what} is in the image, not just \textit{what the pixel values are}.
\end{itemize}

The \textbf{temporal smoothing buffers} we implemented (10-frame majority vote
for mode selection, 5-frame average for atmospheric light) were the most
practically impactful fix we made — turning a jittery, unusable prototype into
a stable, smooth-looking system.

The \textbf{ego-vehicle masking trick} (zeroing the top 60\% of the segmentation
output) solved a frustrating false-positive problem with a single line of code,
demonstrating that sometimes the most elegant engineering solutions are also the
simplest.

Quantitatively, the system achieved sustained throughput above 20 FPS during
all tested conditions on the RTX 3050, with consistent Contrast Gain above 40\%
and Visibility Score recovery above 65\% safe-zone coverage in foggy sequences.

As a future scope, the weather classifier could be fine-tuned on a real dashcam
weather dataset (the model currently uses ImageNet weights transferred to this
domain with zero additional training). The transmission map estimation could
also be refined using a Guided Image Filter instead of a simple Gaussian blur,
which would sharpen object edges in the dehazed output.

Overall, this project gave us hands-on experience in building a complete
computer vision pipeline from scratch — from raw camera input all the way to a
live annotated dashboard — which we believe is much closer to real-world
engineering practice than any isolated lab exercise.

% ============================================================
%  REFERENCES
% ============================================================
\newpage
\begin{thebibliography}{9}

\bibitem{he2011}
K.~He, J.~Sun, and X.~Tang,
\textit{Single Image Haze Removal Using Dark Channel Prior},
IEEE Transactions on Pattern Analysis and Machine Intelligence,
vol.~33, no.~12, pp.~2341--2353, 2011.

\bibitem{sandler2018}
M.~Sandler, A.~Howard \textit{et al.},
\textit{MobileNetV2: Inverted Residuals and Linear Bottlenecks},
Proceedings of the IEEE Conference on Computer Vision and Pattern Recognition
(CVPR), 2018.

\bibitem{howard2019}
A.~Howard \textit{et al.},
\textit{Searching for MobileNetV3},
Proceedings of the IEEE International Conference on Computer Vision (ICCV), 2019.

\bibitem{pizer1987}
S.~M.~Pizer \textit{et al.},
\textit{Adaptive Histogram Equalization and Its Variations},
Computer Vision, Graphics, and Image Processing,
vol.~39, no.~3, pp.~355--368, 1987.

\end{thebibliography}

\end{document}
% ============================================================
%  END OF DOCUMENT
% ============================================================